\documentclass[10pt,a4paper]{amsart}
\usepackage[margin=19mm]{geometry}
\usepackage{amsmath,amssymb,mathtools}
\usepackage[scr=rsfs,cal=euler]{mathalfa}
\usepackage{xfrac}
\usepackage[unicode,hidelinks,bookmarks=false]{hyperref}
\newcommand{\ie}{i.e.\ }
\newcommand{\cf}{cf.\ }
\newcommand{\resp}{respectively\ }
\newcommand{\Subsection}{\subsection}
\providecommand{\nobreakdash}{\nobreak\hbox{-}}
\setlength{\emergencystretch}{2em}
\multlinegap0pt
\usepackage{amssymb}
\usepackage{xcolor}
\usepackage[shortlabels]{enumitem}
\usepackage{tikz}
\newtheorem{lemma}{Lemma}
\title{Inverse Scattering for Single Photons in Quantum Optics}
\author{Giovanni Covi, Matti Lassas, Medet Nursultanov, Lauri Oksanen, Mikko Salo, John C. Schotland}
\date{Source: arXiv:2609.03726v1 (CC BY 4.0)}
\begin{document}
\maketitle

\noindent\emph{Source and licence.} Inverse Scattering for Single Photons in Quantum Optics. Version arXiv:2609.03726v1 (CC BY 4.0), distributed by arXiv under the Creative Commons Attribution 4.0 International licence, http://creativecommons.org/licenses/by/4.0/, as recorded in the arXiv licence metadata for this exact version and shown on the abstract page https://arxiv.org/abs/2609.03726v1. Full licence notice: \texttt{licenses/arxiv-2609.03726-CC-BY-4.0.txt}.\par\smallskip

\noindent\textit{Editorial scope.} Counted unit: the article's lemma lem:aux\_sing (source0.tex lines 2426--2455). The article's complete original proof of it is delivered with it (source0.tex lines 2457--2694, one \texttt{proof} environment of 238 lines). No mathematics is written, repaired or re-derived: the statement and the proof are the article's own lines, in the article's own notation, and no retained line is substituted or re-flowed.  No further source-local context is needed: every cross-reference of every retained line resolves inside the two delivered spans.  Nothing was omitted from any retained span, so the package carries no omission record.  No retained line contains a citation, so the wrapper carries no bibliography and no bibliography entry.  No retained line contains a figure environment, an image-inclusion command or drawing code, so no image is delivered and the package declares no asset dependency.  Editorial formatting only, in addition: an AMS \texttt{amsart} preamble that restates the article's own declarations and macros used by the retained lines, namely the environments \texttt{lemma} on separate counters, together with the standard packages those lines need.  The retained environments print in this document as Lemma 1 in order of appearance, whereas the article prints the same items with the numbering of its own full document.  The complete native source of the article as distributed in the arXiv e-print for this exact version is bundled unchanged as \texttt{sources/209312/source0.tex}.
\begin{lemma}\label{lem:aux_sing}
Let $k> 1$.  Define
\begin{equation*}
    \sigma_k(\eta):=\eta_1+\frac{|\eta|^2}{2k},
  \qquad
  a_k(\eta):=\frac{|ke_1+\eta|+k}{2k},
\end{equation*}
and for $\theta,\vartheta\in[0,1]$ set
\begin{equation}\label{eq:aux_interp}
  \sigma_k^\vartheta(\eta):=\eta_1+\frac{\vartheta|\eta|^2}{2k},
  \qquad
  c_\theta(\eta):=\theta a_k(\eta)+(1-\theta).
\end{equation}
For $\varepsilon>0$, set
\begin{equation}\label{eq:aux_X}
  X_{\theta,\vartheta}(\eta)
  :=\sigma_k^\vartheta(\eta)-i\varepsilon c_\theta(\eta).
\end{equation}
Then there exist constant $C>0$, depending only on $n$, such that for all $k>1$,
$\theta,\vartheta\in[0,1]$, $\varepsilon\in(0,1]$,
and all $\psi\in\mathscr{S}(\mathbb{R}^n)$ supported in $\{|\eta|\leq k/2\}$:
\begin{align}
  \left|\int_{\mathbb{R}^n}\frac{\psi(\eta)}{X_{\theta,\vartheta}(\eta)}\,d\eta\right|
  &\leq C\|\psi\|_{\mathscr{S}_{n+4}},
  \label{eq:aux_inv1}\\
  \left|\int_{\mathbb{R}^n}\frac{\psi(\eta)}{X_{\theta,\vartheta}(\eta)^2}\,d\eta\right|
  &\leq C\|\psi\|_{\mathscr{S}_{n+5}}.
  \label{eq:aux_inv2}
\end{align}
\end{lemma}

\begin{proof}
We write $X=X_{\theta,\vartheta}$. Next, we list certain properties of functions $\sigma_k^\vartheta$, $c_\theta$, and $X$ on $\{|\eta|\leq k/2\}$.


\textbf{Properties of $\sigma_k^\vartheta$.}
We compute
\begin{equation*}
    \partial_{\eta_1}\sigma_k^\vartheta = 1+ \frac{\vartheta\eta_1}{k},
\end{equation*}
so
\begin{equation}\label{eq:aux_sigma_bounds}
  \frac{1}{2}\leq \partial_{\eta_1}\sigma_k^\vartheta\le\frac{3}{2}
  \qquad
  \text{on }\{|\eta_1|\leq k/2\}.
\end{equation}
Moreover,
\begin{equation*}
    \partial_{\eta_i}\partial_{\eta_j}\sigma_k^\vartheta = \frac{\vartheta\delta_{ij}}{k} \leq \frac{1}{k}
\end{equation*}
and all derivatives of order greater than $2$ are zero. Hence,
\begin{equation}\label{eq:aux_sigma_allderivs}
  |\partial_\eta^\nu\sigma_k^\vartheta|\leq C_\nu
  \qquad\text{on }\{|\eta|\leq k/2\},\quad\text{for }|\nu|\geq 1,
\end{equation}
uniformly in $k> 1$ and $\vartheta\in[0,1]$.

For $\vartheta=0$, $\sigma_k^0=\eta_1$ has the unique zero $\eta_1^*=0$.
For $\vartheta>0$, the equation $\sigma_k^\vartheta(\eta_1,\eta')=0$
is quadratic in $\eta_1$ with two roots:
\begin{equation*}
    \eta_1^*=\frac{k}{\vartheta}\Bigl(-1+\sqrt{1-\tfrac{\vartheta^2|\eta'|^2}{k^2}}\Bigr),
  \qquad
  \eta_1^{**}=\frac{k}{\vartheta}\Bigl(-1-\sqrt{1-\tfrac{\vartheta^2|\eta'|^2}{k^2}}\Bigr).
\end{equation*}

Since $|\eta_1^{**}|\geq k/\vartheta\geq k$, we obtain $\eta_1^{**} \notin \{|\eta_1|\leq k/2\}$. For $\eta' \in \{|\eta'|\leq k/2\}$, we estimate
\begin{equation}\label{eq:aux_eta1star}
    |\eta_1^*|=\frac{k}{\vartheta}\Bigl(1-\sqrt{1-\tfrac{\vartheta^2|\eta'|^2}{k^2}}\Bigr) \leq \frac{\frac{\vartheta|\eta'|^2}{k}}{1+\sqrt{1-\tfrac{\vartheta^2|\eta'|^2}{k^2}}} \leq \frac{|\eta'|^2}{k} \leq \frac{k}{4} < \frac{k}{2}.
\end{equation}
Hence, $\eta_1^* \in \{|\eta_1|\leq k/2\}$. Using the mean value theorem and \eqref{eq:aux_sigma_bounds}, we obtain the two-sided bound:
\begin{equation}\label{eq:aux_sigma_two_sided}
  \tfrac{1}{2}|\eta_1-\eta_1^*|
  \le|\sigma_k^\vartheta(\eta_1,\eta')|
  \le\tfrac{3}{2}|\eta_1-\eta_1^*|,
  \qquad
  \text{for } \eta \in \left\{|\eta|<k/2\right\}.
\end{equation}
Note that in general $(\eta_1^*,\eta')$ is not in $\left\{|\eta|<k/2\right\}$.

\textbf{Properties of $c_\theta$.}
Note that 
\begin{equation*}
    \frac{k}{2}\leq |ke_1+\eta|\leq \frac{3k}{2},
    \qquad
    \text{on }
    \{|\eta|\leq k/2\}
\end{equation*}
Therefore, 
\begin{equation}\label{eq:aux_c_bounds}
  \frac{3}{4}\leq c_\theta\le\frac{5}{4}
  \qquad\text{on }\{|\eta|\leq k/2\}.
\end{equation}
Writing $\xi=ke_1+\eta$ with $|\xi|\geq k/2$, we have
$\partial_{\eta_j}a_k=\xi_j/(2k|\xi|)$, and each further $\eta$-derivative
brings a factor of $|\xi|^{-1}\leq 2/k$.
Hence
\begin{equation}\label{eq:aux_a_allderivs}
    |\partial_\eta^\nu a_k|\leq C_\nu k^{-|\nu|},
    \qquad
    \text{for }
    |\nu|\geq 1
\end{equation}
and therefore
\begin{equation}\label{eq:aux_c_allderivs}
  |\partial_\eta^\nu c_\theta|\leq C_\nu k^{-|\nu|}\leq C_\nu
  \qquad\text{on }\{|\eta|\leq k/2\},\quad|\nu|\geq 1,
\end{equation}
uniformly in $k> 1$ and $\theta\in[0,1]$.

\textbf{Properties of $X$.}
Due to \eqref{eq:aux_c_bounds},
\begin{equation}\label{eq:aux_Xnonzero}
  |X|= |\sigma_k^\vartheta-i\varepsilon c_\theta| \geq \varepsilon c_\theta\geq \frac{3}{4}\varepsilon>0,
  \qquad
  \text{on } \{|\eta|\leq k/2\},
\end{equation}
and hence, $X$ does not vanish on $\{|\eta|\leq k/2\}$.
Also $\operatorname{Im} X=-\varepsilon c_\theta<0$, so $X$ lies in the open
lower half-plane. By \eqref{eq:aux_sigma_bounds}
\begin{equation}\label{eq:aux_dX_lower}
    |\partial_{\eta_1}X|
    =\sqrt{(\partial_{\eta_1}\sigma_k^\vartheta)^2
    +\varepsilon^2(\partial_{\eta_1}c_\theta)^2}
    \geq |\partial_{\eta_1}\sigma_k^\vartheta|
    \geq\frac{1}{2}.
\end{equation}
Estimates \eqref{eq:aux_sigma_allderivs} and \eqref{eq:aux_c_allderivs} imply that
\begin{equation}\label{eq:aux_X_allderivs}
  |\partial_\eta^\nu X|\leq C_\nu
  \qquad\text{on }\{|\eta|\leq k/2\},\quad\text{for }|\nu|\geq 1,
\end{equation}
uniformly in $k> 1$, $\theta,\vartheta\in[0,1]$, and $\varepsilon\in(0,1]$.

Combining the last two estimates gives 
\begin{equation}\label{eq:aux_inv_dX_allderivs}
  \left|\partial_\eta^\nu\left[\frac{1}{\partial_{\eta_1}X}\right]\right|\leq C_\nu
  \qquad\text{on }\{|\eta|\leq k/2\},\quad\text{for all }\nu.
\end{equation}


\textbf{Proof of \eqref{eq:aux_inv1}.}
As we established, $X$ and $\partial_{\eta_1} X$ do not vanish on $\operatorname{supp} (\psi) \subset \{|\eta|\leq k/2\}$. Therefore, by using
\begin{equation*}
    \frac{1}{X} = \frac{\partial_{\eta_1}[\log X]}{\partial_{\eta_1} X}
\end{equation*}
and integrating by parts, we obtain
\begin{equation}\label{eq:aux_IBP}
  \int_{\mathbb{R}^n}\frac{\psi}{X}\,d\eta
  = -\int_{\mathbb{R}^n}\log X(\eta)\;\partial_{\eta_1}\!\left[
      \frac{\psi(\eta)}{\partial_{\eta_1}X(\eta)}
    \right]d\eta.
\end{equation}
Here $\log X$ is the principal-branch logarithm, well-defined since
$\operatorname{Im} X<0$ (so $\arg X\in(-\pi,0)$). Therefore,
\begin{equation}\label{eq:logX}
    |\log X| \leq |\log |X|| + |\arg X| \leq |\log |X|| + \pi.
\end{equation}

By the definition of $X$ and \eqref{eq:aux_c_bounds},
\begin{equation*}
    |\sigma_k^\vartheta| \leq |X| \leq |\sigma_k^\vartheta|  + \frac{5}{4}.
\end{equation*}
For $|X|\geq 1$, we estimate
\begin{equation*}
    \log |X| \leq \log \left(|\sigma_k^\vartheta|  + \frac{5}{4}\right) \leq 
    \begin{cases}
        \log\left(\frac{9}{4}|\sigma_k^\vartheta| \right) & |\sigma_k^\vartheta| \geq 1,\\
        \log\left(\frac{9}{4} \right) & |\sigma_k^\vartheta| \leq 1.\
    \end{cases}
\end{equation*}
For $|X|\leq 1$,
\begin{equation*}
    |\log |X|| = - \log |X|\leq - \log |\sigma_k^\vartheta| = |\log |\sigma_k^\vartheta||.
\end{equation*}
In both cases,
\begin{equation*}
    |\log |X|| \leq |\log |\sigma_k^\vartheta|| + \log\left(\frac{9}{4} \right).
\end{equation*}
Combining this with \eqref{eq:aux_sigma_two_sided} and \eqref{eq:logX}, we obtain
\begin{equation}\label{eq:aux_logX}
  |\log X(\eta)|\le|\log|\eta_1-\eta_1^*(\eta')||+C
\end{equation}
for some absolute constant $C$.

Next, we estimate the second term in the integral \eqref{eq:aux_IBP}. Set $\varphi(\eta):=\partial_{\eta_1}[\psi/\partial_{\eta_1}X]$.
Since $\psi$ is smooth and supported in $\{|\eta|\leq k/2\}$, \eqref{eq:aux_inv_dX_allderivs} implies that
\begin{equation}\label{eq:aux_phi_bound}
  \|\varphi\|_{\mathscr{S}_m}\leq C\|\psi\|_{\mathscr{S}_{m+1}}
\end{equation}
for every $m$, with $C=C(m,n)$.

Combining \eqref{eq:aux_IBP} and \eqref{eq:aux_logX} gives
\begin{equation*}
    \left|\int_{\mathbb{R}^n}\frac{\psi}{X} d\eta\right|
    \leq C\int_{\mathbb{R}^n}\left(1+|\log|\eta_1-\eta_1^*(\eta')||\right) |\varphi(\eta)| d\eta.
\end{equation*}
We split the inner $\eta_1$-integral at $|\eta_1-\eta_1^*|=1$:

\begin{multline}\label{eq:spl_int}
    \left|\int_{\mathbb{R}^n}\frac{\psi}{X} d\eta\right|
  \leq 
  C\int_{\mathbb{R}^n} |\varphi(\eta)| d\eta
  + C \int_{\mathbb{R}^{n-1}}\int_{|\eta_1 - \eta_1^*(\eta')| \leq 1} |\log|\eta_1-\eta_1^*(\eta')| |\varphi(\eta)| d\eta\\
  + C \int_{\mathbb{R}^{n-1}}\int_{|\eta_1 - \eta_1^*(\eta')| \geq 1} |\log|\eta_1-\eta_1^*(\eta')| |\varphi(\eta)| d\eta.
\end{multline}
Since $\varphi \in C_c^\infty(\mathbb{R}^n)$,
\begin{equation*}
    |\varphi(\eta)|\le\|\varphi\|_{\mathscr{S}_{p+1}}(1+|\eta|)^{-p}
    \qquad
    \text{for any } p\in \mathbb{N}.
\end{equation*}
Choosing $p = n + 2$, we estimate the first term of the right-hand side of \eqref{eq:spl_int}
\begin{equation*}
    \int_{\mathbb{R}^n} |\varphi(\eta)| d\eta \leq C \|\varphi\|_{\mathscr{S}_{n+3}}.
\end{equation*}
To estimate the second term, we note that
\begin{equation*}
    \int_{|\eta_1-\eta_1^*|\leq 1}\frac{|\log|\eta_1-\eta_1^*||}{(1+|\eta|)^p} d\eta_1
    \leq \frac{1}{(1+|\eta'|)^p}\int_{|s|\leq 1}|\log|s|| ds
    = \frac{2}{(1+|\eta'|)^p}.
\end{equation*}
Hence,
\begin{equation*}
    \int_{\mathbb{R}^{n-1}}\int_{|\eta_1 - \eta_1^*(\eta')| \leq 1} |\log|\eta_1-\eta_1^*(\eta')| |\varphi(\eta)| d\eta \leq  C \|\varphi\|_{\mathscr{S}_{n+3}}.
\end{equation*}
To estimate the last term of the of the right-hand side of \eqref{eq:spl_int}, by using \eqref{eq:aux_eta1star} and $\log|t|\le|t|$ for $|t|>1$, we obtain
\begin{equation*}
    \frac{|\log|\eta_1-\eta_1^*||}{(1+|\eta|)^p}
    \leq 
    \frac{|\eta_1-\eta_1^*|}{(1+|\eta|)^p}
    \leq
    \frac{|\eta_1|+|\eta_1^*|}{(1+|\eta|)^p}
    \leq
    \frac{C(1+|\eta|)}{(1+|\eta|)^p}=\frac{C}{(1+|\eta|)^{p-1}}.
\end{equation*}
Therefore,
\begin{equation*}
    \int_{\mathbb{R}^{n-1}}\int_{|\eta_1 - \eta_1^*(\eta')| \geq 1} |\log|\eta_1-\eta_1^*(\eta')| |\varphi(\eta)| d\eta \leq  C \|\varphi\|_{\mathscr{S}_{n+3}}.
\end{equation*}
By combining estimates for the three terms of the right-hand side of \eqref{eq:spl_int}, and then using \eqref{eq:aux_phi_bound}, 
\begin{equation*}
    \left|\int_{\mathbb{R}^n}\frac{\psi}{X}d\eta\right|
    \leq C\|\varphi\|_{\mathscr{S}_{n+3}}
    \leq C\|\psi\|_{\mathscr{S}_{n+4}}
\end{equation*}
This proves \eqref{eq:aux_inv1}.

\textbf{Proof of \eqref{eq:aux_inv2}.}
Since $X$ and $\partial_{\eta_1} X$ do not vanish on $\operatorname{supp} (\psi) \subset \{|\eta|\leq k/2\}$, we can write 
\begin{equation*}
    \frac{1}{X^2} = - \frac{1}{\partial_{\eta_1}X} \partial_{\eta_1} \left[\frac{1}{X}\right]
\end{equation*}
on $\{|\eta|\leq k/2\}$. Then, the integration by parts gives
\begin{equation*}
    \int_{\mathbb{R}^n}\frac{\psi}{X^2}d\eta
    = \int_{\mathbb{R}^n}\frac{1}{X}
    \partial_{\eta_1}\left[\frac{\psi}{\partial_{\eta_1}X}\right]d\eta = \int_{\mathbb{R}^n}\frac{\varphi}{X}
    d\eta.
\end{equation*}
 Recall that $\varphi$ is a smooth function supported in $\{|\eta|\leq k/2\}$. Therefore, applying \eqref{eq:aux_inv1} and \eqref{eq:aux_phi_bound}, we obtain 
\begin{equation*}
    \left|\int_{\mathbb{R}^n}\frac{\psi}{X^2} d\eta\right| 
    = \left|\int_{\mathbb{R}^n}\frac{\varphi}{X} d\eta\right| 
    \leq C\|\varphi\|_{\mathscr{S}_{n+4}} 
    \leq C\|\psi\|_{\mathscr{S}_{n+5}}.
\end{equation*}
This gives \eqref{eq:aux_inv2}.
\end{proof}

\end{document}
